% รายงานความก้าวหน้า รอบ 9 เดือน — กลุ่มแผนงานใต้ร่มพระบารมี มทร.ล้านนา ปีงบประมาณ 2569
% ตัวเลขทั้งหมดมาจาก: ไฟล์งบประมาณ 14/7/2569 · ไฟล์จัดสรรแผนงานและตัวชี้วัด 20/7/2569 · ฐานข้อมูลติดตามโครงการ ณ 22/7/2569
% build: xelatex -interaction=nonstopmode RPF2569_progress_report.tex ×2
\documentclass[12pt,a4paper]{article}

\usepackage[a4paper,top=2.2cm,bottom=2.2cm,left=2.3cm,right=2cm,headheight=15pt]{geometry}
\usepackage{polyglossia}
\setdefaultlanguage{thai}
\setotherlanguage{english}
\IfFontExistsTF{TH Sarabun New}{%
  \setmainfont[Ligatures=TeX]{TH Sarabun New}%
  \newfontfamily\thaifont[Script=Thai,Ligatures=TeX]{TH Sarabun New}%
}{%
  \setmainfont[Ligatures=TeX]{TH SarabunPSK}%
  \newfontfamily\thaifont[Script=Thai,Ligatures=TeX]{TH SarabunPSK}%
}
\XeTeXlinebreaklocale "th"
\XeTeXlinebreakskip = 0pt plus 1pt
\setmonofont[Scale=0.92]{Menlo}

\usepackage{amsmath}
\usepackage{xcolor}
\usepackage{colortbl}
\usepackage{longtable}
\usepackage{booktabs}
\usepackage{array}
\usepackage{tikz}
\usepackage{fancyhdr}
\usepackage{titlesec}
\usepackage{enumitem}
\usepackage[hidelinks]{hyperref}

\definecolor{rpfnavy}{RGB}{25,58,110}
\definecolor{rpfgold}{RGB}{191,146,42}
\definecolor{rpfteal}{RGB}{15,118,110}
\definecolor{rpfink}{RGB}{31,45,58}
\definecolor{rpfgreen}{RGB}{22,132,74}
\definecolor{rpfamber}{RGB}{202,138,4}
\definecolor{rpfred}{RGB}{185,45,45}

\pagestyle{fancy}\fancyhf{}
\fancyhead[L]{\footnotesize รายงานความก้าวหน้า รอบ 9 เดือน $\cdot$ กลุ่มแผนงานใต้ร่มพระบารมี}
\fancyhead[R]{\footnotesize ปีงบประมาณ 2569}
\fancyfoot[C]{\footnotesize\thepage}
\renewcommand{\headrulewidth}{0.3pt}
\setlength{\parindent}{0pt}
\setlength{\parskip}{4pt}

\titleformat{\section}{\Large\bfseries\color{rpfnavy}}{\thesection.}{8pt}{}
\titleformat{\subsection}{\large\bfseries\color{rpfteal}}{\thesubsection}{8pt}{}
\titlespacing{\section}{0pt}{14pt}{6pt}

% กล่องเน้นข้อความ
\newcommand{\keybox}[2]{%
  \begin{center}\setlength{\fboxsep}{8pt}%
  \fcolorbox{#1}{#1!8}{\begin{minipage}{0.94\linewidth}\color{rpfink}#2\end{minipage}}%
  \end{center}}

\begin{document}

% ================= ปก =================
\begin{titlepage}
\thispagestyle{empty}
\begin{tikzpicture}[remember picture,overlay]
  \fill[rpfnavy] (current page.north west) rectangle ([yshift=-11.5cm]current page.north east);
  \fill[rpfgold] ([yshift=-11.5cm]current page.north west) rectangle ([yshift=-11.9cm]current page.north east);
  % ลายภูเขา
  \fill[white,opacity=0.07]
    ([xshift=0cm,yshift=-9.6cm]current page.north west)
    -- ++(4.5,2.1) -- ++(3.2,-1.4) -- ++(4.0,2.4) -- ++(3.5,-1.7) -- ++(6.0,2.6)
    |- ([yshift=-11.5cm]current page.north east) -- cycle;
\end{tikzpicture}
\vspace*{2.1cm}
{\color{white}\fontsize{15}{18}\selectfont รายงานความก้าวหน้าการดำเนินงาน รอบ 9 เดือน}\par\vspace{10pt}
{\color{white}\fontsize{34}{40}\selectfont\bfseries กลุ่มแผนงานใต้ร่มพระบารมี}\par\vspace{8pt}
{\color{rpfgold}\fontsize{22}{26}\selectfont\bfseries มหาวิทยาลัยเทคโนโลยีราชมงคลล้านนา}\par\vspace{14pt}
{\color{white}\fontsize{17}{20}\selectfont ปีงบประมาณ พ.ศ. 2569}\par\vspace{4pt}
{\color{white!80}\fontsize{13}{16}\selectfont ข้อมูล ณ วันที่ 22 กรกฎาคม 2569 $\cdot$ 1 ตุลาคม 2568 -- 22 กรกฎาคม 2569}\par
\vspace{3.4cm}

\begin{center}
\setlength{\tabcolsep}{10pt}\renewcommand{\arraystretch}{1.5}
\begin{tabular}{>{\centering\arraybackslash}p{3.1cm}>{\centering\arraybackslash}p{3.1cm}>{\centering\arraybackslash}p{3.1cm}>{\centering\arraybackslash}p{3.1cm}}
\rowcolor{rpfnavy!8}
{\color{rpfnavy}\Large\bfseries 61} & {\color{rpfnavy}\Large\bfseries 3} & {\color{rpfnavy}\Large\bfseries 11} & {\color{rpfnavy}\Large\bfseries 48}\\
\rowcolor{rpfnavy!8}
{\footnotesize โครงการย่อย} & {\footnotesize โครงการหลัก ง8} & {\footnotesize หน่วยงานร่วม} & {\footnotesize พื้นที่ดำเนินงาน}\\
\end{tabular}
\par\vspace{16pt}
\setlength{\tabcolsep}{12pt}
\begin{tabular}{>{\centering\arraybackslash}p{5.1cm}>{\centering\arraybackslash}p{5.1cm}}
\rowcolor{rpfgold!12}
{\color{rpfnavy}\Large\bfseries 7,986,583} & {\color{rpfnavy}\Large\bfseries 4,017,804}\\
\rowcolor{rpfgold!12}
{\footnotesize งบประมาณที่ได้รับจัดสรร (บาท)} & {\footnotesize เบิกจ่ายแล้ว (บาท) คิดเป็น 50\%}\\
\end{tabular}
\end{center}
\vfill
{\footnotesize\color{rpfink}
แหล่งข้อมูล: ไฟล์งบประมาณกลุ่มแผนงานใต้ร่มพระบารมี ปีงบประมาณ 2569 (ฉบับ 14 กรกฎาคม 2569) $\cdot$
ไฟล์จัดสรรแผนงาน งบประมาณ และตัวชี้วัด (ปรับปรุง 20 กรกฎาคม 2569) $\cdot$
ฐานข้อมูลระบบติดตามโครงการ (61 โครงการ 268 กิจกรรม)}
\end{titlepage}

% ================= สารบัญ =================
\tableofcontents
\clearpage

% ================= บทสรุปผู้บริหาร =================
\section{บทสรุปผู้บริหาร}

กลุ่มแผนงานใต้ร่มพระบารมี มหาวิทยาลัยเทคโนโลยีราชมงคลล้านนา ดำเนินงานปีงบประมาณ 2569 ภายใต้
\textbf{3 โครงการหลักตามแบบ ง8} รวม \textbf{61 โครงการย่อย} งบประมาณที่ได้รับจัดสรร \textbf{7,986,583 บาท}
โดยมีหน่วยงานร่วมดำเนินการ 11 หน่วยงาน หัวหน้าโครงการ 42 ท่าน ครอบคลุมพื้นที่ดำเนินงาน 48 พื้นที่
ในศูนย์พัฒนาโครงการหลวง สถานีเกษตรหลวง โครงการตามพระราชดำริ และชุมชนบนพื้นที่สูง

\keybox{rpfgreen}{%
\textbf{\color{rpfgreen}ผลสำเร็จที่โดดเด่น}\par
ตัวชี้วัดที่ 39 --- จำนวนองค์ความรู้ เทคโนโลยี และนวัตกรรมที่นำไปใช้ในโครงการหลวง โครงการตามพระราชดำริ
ซึ่งเป็น\textbf{เป้าหมายของกลุ่มแผนงานใต้ร่มพระบารมีโดยตรง} --- รับมาดำเนินการ \textbf{62 องค์ความรู้ จากเป้า 60}
คิดเป็น \textbf{103\%} กระจายใน 37 โครงการ ; โอนงบประมาณถึงผู้รับผิดชอบแล้ว 98\% ;
อนุมัติโครงการแล้ว 59 จาก 61 โครงการ}

\keybox{rpfamber}{%
\textbf{\color{rpfamber}ประเด็นที่ต้องเร่งรัด}\par
การเบิกจ่ายอยู่ที่ \textbf{4,017,804 บาท คิดเป็น 50\%} ขณะที่ระยะเวลาปีงบประมาณผ่านไปแล้ว \textbf{81\%}
เท่ากับ\textbf{ช้ากว่าเวลาประมาณ 31 จุดร้อยละ} คงเหลือที่ต้องเบิกจ่าย 3,968,779 บาท ภายในเวลาประมาณ 70 วัน
โดยมีภาระผูกพันใบสั่งซื้อรองรับแล้วเพียง 730,948 บาท}

\keybox{rpfred}{%
\textbf{\color{rpfred}ความเสี่ยงเชิงระบบ}\par
(1) ตัวชี้วัดเชิงเครือข่ายและเชิงพาณิชย์ยังต่ำกว่า 20\% ได้แก่ ตัวชี้วัดที่ 10 (18\%) ที่ 17 (13\%) และที่ 38 (8\%) ;
(2) \textbf{ยังไม่มีการรายงานผลรายกิจกรรมเข้าระบบติดตาม (0 จาก 268 กิจกรรม)}
ทำให้รายงานตัวชี้วัดในขณะนี้สะท้อน\textbf{ค่าที่รับมาดำเนินการ (commitment)} ไม่ใช่ผลสำเร็จที่ตรวจรับแล้ว}

\subsection*{ข้อเสนอต่อผู้บริหาร}
\begin{enumerate}[leftmargin=1.6em,itemsep=1pt]
  \item เร่งรัดการเบิกจ่ายรายโครงการที่ต่ำกว่า 30\% เป็นรายสัปดาห์ โดยเฉพาะ มทร.ล้านนา ตาก (3 โครงการ 250,000 บาท) และคณะบริหารธุรกิจและศิลปศาสตร์ (1 โครงการ 40,000 บาท) ที่ยังไม่มีการเบิกจ่าย
  \item กำหนดให้ผู้รับผิดชอบโครงการรายงานผลรายกิจกรรมพร้อมหลักฐานเข้าระบบติดตามให้ครบภายในกลางเดือนกันยายน 2569
  \item เร่งยื่นจดทะเบียนทรัพย์สินทางปัญญาจากผลงานที่พร้อม และจัดทำบันทึกความร่วมมือกับหน่วยงานภายนอกและสถานประกอบการ เพื่อปิดช่องว่างตัวชี้วัดที่ 10 17 และ 38
  \item ทบทวนการออกแบบตัวชี้วัดปีงบประมาณ 2570 ให้สอดคล้องกับลักษณะงานพัฒนาชุมชนตั้งแต่ต้นปี
\end{enumerate}

\clearpage
% ================= 1 บทนำ =================
\section{บทนำ}
\subsection{ที่มาและความสำคัญ}
มหาวิทยาลัยเทคโนโลยีราชมงคลล้านนา มีภารกิจสนับสนุนการดำเนินงานโครงการหลวงและโครงการอันเนื่องมาจากพระราชดำริ
บนพื้นที่สูงภาคเหนือมาอย่างต่อเนื่อง เพื่อ\textbf{สืบสาน รักษา และต่อยอด} พระราชปณิธานในการพัฒนาคุณภาพชีวิต
ของประชาชนบนพื้นที่สูง

ในปีงบประมาณ 2569 มหาวิทยาลัยได้รวมงานที่เคยกระจัดกระจายตามคณะและวิทยาเขต เข้าเป็น
\textbf{กลุ่มแผนงานใต้ร่มพระบารมี} แผนงานเดียวที่มีเป้าหมาย ตัวชี้วัด และระบบติดตามร่วมกัน
ครอบคลุม 5 ยุทธศาสตร์ 7 แผนงาน 3 โครงการหลักตามแบบ ง8 และ 61 โครงการย่อย

\subsection{วัตถุประสงค์ของแผนงาน}
\begin{enumerate}[leftmargin=1.6em,itemsep=1pt]
  \item วิจัยและพัฒนาเทคโนโลยีและนวัตกรรมที่ใช้ได้จริงในพื้นที่สูงและพื้นที่โครงการหลวง
  \item ยกระดับคุณภาพชีวิตและเสริมสร้างความเข้มแข็งของชุมชนตามศาสตร์พระราชาและหลักปรัชญาเศรษฐกิจพอเพียง
  \item พัฒนากำลังคน สร้างอาชีพ และลดความเหลื่อมล้ำบนพื้นที่สูง
  \item สร้างกลไกการบริหารจัดการ การติดตามประเมินผล และเครือข่ายความร่วมมือที่ยั่งยืน
\end{enumerate}

\subsection{ขอบเขตและวิธีการจัดทำรายงาน}
รายงานฉบับนี้เป็น\textbf{รายงานความก้าวหน้ารอบ 9 เดือน} (1 ตุลาคม 2568 -- 22 กรกฎาคม 2569)
\textbf{มิใช่รายงานผลการดำเนินงานประจำปีฉบับสมบูรณ์} เนื่องจากปีงบประมาณยังไม่สิ้นสุด
ข้อมูลทั้งหมดตรวจสอบตรงจากสามแหล่ง คือ ไฟล์งบประมาณฉบับ 14 กรกฎาคม 2569
ไฟล์จัดสรรแผนงานและตัวชี้วัดฉบับปรับปรุง 20 กรกฎาคม 2569 และฐานข้อมูลระบบติดตามโครงการ

\keybox{rpfnavy}{\textbf{ข้อจำกัดของข้อมูล} --- ตัวเลขตัวชี้วัดในรายงานนี้คือ\textbf{ค่าเป้าหมายที่แต่ละโครงการรับมาดำเนินการ}
เนื่องจากยังไม่มีการบันทึกผลสำเร็จรายกิจกรรมเข้าระบบติดตาม (0 จาก 268 กิจกรรม)
ผลสำเร็จที่ผ่านการตรวจรับจะปรากฏในรายงานฉบับสมบูรณ์เมื่อสิ้นปีงบประมาณ}

\clearpage
% ================= 2 กรอบการดำเนินงาน =================
\section{กรอบการดำเนินงาน}
\subsection{โครงสร้าง 3 โครงการหลักตามแบบ ง8}

\begin{center}\small
\setlength{\tabcolsep}{5pt}\renewcommand{\arraystretch}{1.35}
\begin{tabular}{p{6.6cm}rrrr}
\toprule
\rowcolor{rpfnavy!12}
\textbf{โครงการหลัก (ง8)} & \textbf{โครงการ} & \textbf{งบจัดสรร} & \textbf{เบิกจ่าย} & \textbf{ร้อยละ}\\
\midrule
ง8-1 ผลักดันเทคโนโลยี นวัตกรรมสู่ชุมชน ตามเป้าหมายการพัฒนาอย่างยั่งยืน & 16 & 1,988,000 & 943,416 & 48\%\\
ง8-2 ขับเคลื่อนกลไกการพัฒนาองค์ความรู้เพื่อยกระดับคุณภาพชีวิต & 14 & 1,999,010 & 1,276,760 & 66\%\\
ง8-3 พัฒนากำลังคน สร้างอาชีพ ลดความเหลื่อมล้ำ บนพื้นที่สูง & 31 & 3,999,573 & 1,797,628 & 46\%\\
\midrule
\rowcolor{rpfgold!15}
\textbf{รวมทั้งแผนงาน} & \textbf{61} & \textbf{7,986,583} & \textbf{4,017,804} & \textbf{50\%}\\
\bottomrule
\end{tabular}
\end{center}

\subsection{ยุทธศาสตร์และแผนงานที่เกี่ยวข้อง}
แผนงานทั้ง 61 โครงการอยู่ภายใต้ 5 ยุทธศาสตร์ และ 7 แผนงานย่อย ได้แก่
แผนงานที่ 1.1 วิจัยและพัฒนาเทคโนโลยีนวัตกรรมสร้างรายได้บนพื้นที่สูง $\cdot$
1.2 วิจัยและพัฒนาเพื่อแก้ปัญหาเชิงพื้นที่อย่างมีส่วนร่วม $\cdot$
2.1 วิจัยและพัฒนาเพื่อส่งเสริมอาชีพในภาคเกษตรและนอกภาคเกษตร $\cdot$
2.2 วิจัยและพัฒนาภายใต้มาตรฐานอาหารปลอดภัยและเป็นมิตรต่อสิ่งแวดล้อม $\cdot$
3.1 ยกระดับคุณภาพชีวิตและความเข้มแข็งของชุมชนตามศาสตร์พระราชา $\cdot$
4.1 สนับสนุนการพัฒนาโครงสร้างพื้นฐานเพื่อเพิ่มประสิทธิภาพการดำเนินงาน $\cdot$
5.1 ขับเคลื่อนกลไกการปฏิบัติงานอย่างมีส่วนร่วมสู่การพัฒนาที่ยั่งยืน

\subsection{ทฤษฎีและกรอบแนวคิดที่นำมาใช้}
การดำเนินงานทั้งแผนงานวางอยู่บนฐานคิดที่ปรากฏเป็นชื่อแผนงานและโครงการจริงในปี 2569 ดังนี้

\begin{center}\small
\setlength{\tabcolsep}{6pt}\renewcommand{\arraystretch}{1.3}
\begin{tabular}{p{4.6cm}p{10.4cm}}
\toprule
\rowcolor{rpfnavy!12}
\textbf{กรอบแนวคิด} & \textbf{การนำมาใช้จริงในแผนงาน}\\
\midrule
ศาสตร์พระราชาและปรัชญาเศรษฐกิจพอเพียง & แผนงานที่ 3.1 ยกระดับคุณภาพชีวิตชุมชนอยู่ดีมีสุขตามศาสตร์พระราชา\\
เกษตรทฤษฎีใหม่ และ โคก หนอง นา & โครงการพัฒนาพื้นที่เกษตรทฤษฎีใหม่ตามหลักปรัชญาเศรษฐกิจพอเพียง (555,000 บาท) และโครงการพัฒนาพื้นที่ต้นแบบ โคก หนอง นา\\
เข้าใจ เข้าถึง พัฒนา & โครงการยกระดับมาตรฐานผลิตภัณฑ์ผ้าใยกัญชงวิสาหกิจชุมชนดาวม่างสู่มาตรฐาน มผช.\\
การวิจัยแบบมีส่วนร่วม (Participatory Research) & แผนงานที่ 1.2 พัฒนากระบวนการวิจัยแบบมีส่วนร่วมในระดับพื้นที่ (600,000 บาท)\\
Re-skill Up-skill New-skill & โครงการพัฒนาฐานข้อมูลด้านวิศวกรรมและยกระดับทักษะอาชีพในพื้นที่โครงการหลวง\\
เครือข่ายความร่วมมือ (MOU / MOA) & โครงการจัดทำแผนความร่วมมือและสร้างกลไกพัฒนาเครือข่ายเชิงพื้นที่\\
Logic Model (Output--Outcome--Impact) & กรอบวิเคราะห์ผลการดำเนินงานราย ง8 ในบทที่ 4 ของรายงานฉบับนี้\\
\bottomrule
\end{tabular}
\end{center}

\clearpage
% ================= 3 ผลงบประมาณ =================
\section{ผลการดำเนินงานด้านงบประมาณ}
\subsection{ภาพรวมการใช้จ่ายงบประมาณ}

\begin{center}\small
\setlength{\tabcolsep}{7pt}\renewcommand{\arraystretch}{1.35}
\begin{tabular}{p{5.4cm}rr}
\toprule
\rowcolor{rpfnavy!12}
\textbf{รายการ} & \textbf{จำนวนเงิน (บาท)} & \textbf{ร้อยละ}\\
\midrule
งบประมาณที่ได้รับจัดสรร & 7,986,583 & 100\%\\
โอนงบประมาณให้ผู้รับผิดชอบแล้ว & 7,798,583 & 98\%\\
ผูกพันใบสั่งซื้อ (ยังไม่เบิก) & 730,948 & 9\%\\
\rowcolor{rpfgreen!10}
\textbf{เบิกจ่ายจริง} & \textbf{4,017,804} & \textbf{50\%}\\
เบิกจ่ายรวมภาระผูกพัน & 4,748,752 & 59\%\\
\rowcolor{rpfamber!12}
\textbf{คงเหลือที่ต้องเบิกจ่าย} & \textbf{3,968,779} & \textbf{50\%}\\
\bottomrule
\end{tabular}
\end{center}

\keybox{rpfamber}{\textbf{เทียบกับกรอบเวลา} --- ณ 22 กรกฎาคม 2569 ระยะเวลาปีงบประมาณผ่านไป \textbf{81\%}
ขณะที่เบิกจ่ายได้ \textbf{50\%} ต่างกัน \textbf{31 จุดร้อยละ} เหลือเวลาถึง 30 กันยายน 2569 ประมาณ \textbf{70 วัน}}

\subsection{การเบิกจ่ายรายโครงการหลัก ง8}

\begin{center}\small
\setlength{\tabcolsep}{5pt}\renewcommand{\arraystretch}{1.35}
\begin{tabular}{p{3.5cm}rrrrr}
\toprule
\rowcolor{rpfnavy!12}
\textbf{โครงการหลัก} & \textbf{จัดสรร} & \textbf{โอนแล้ว} & \textbf{ใบสั่งซื้อ} & \textbf{เบิกจ่าย} & \textbf{คงเหลือ}\\
\midrule
ง8-1 ผลักดันเทคโนโลยี & 1,988,000 & 1,950,000 & 173,100 & 943,416 & 1,006,584\\
ง8-2 ขับเคลื่อนองค์ความรู้ & 1,999,010 & 1,948,510 & 114,588 & 1,276,760 & 671,750\\
ง8-3 พัฒนากำลังคน & 3,999,573 & 3,900,073 & 443,260 & 1,797,628 & 2,102,445\\
\midrule
\rowcolor{rpfgold!15}
\textbf{รวม} & \textbf{7,986,583} & \textbf{7,798,583} & \textbf{730,948} & \textbf{4,017,804} & \textbf{3,780,779}\\
\bottomrule
\end{tabular}
\end{center}
{\footnotesize หมายเหตุ: ยอดคงเหลือรายโครงการหลักคำนวณจากยอดโอนงบประมาณ จึงต่างจากยอดคงเหลือรวมที่คำนวณจากงบจัดสรร (3,968,779 บาท)}

\subsection{การเบิกจ่ายรายหน่วยงาน}

\begin{center}\small
\setlength{\tabcolsep}{6pt}\renewcommand{\arraystretch}{1.3}
\begin{tabular}{p{6.6cm}rrrr}
\toprule
\rowcolor{rpfnavy!12}
\textbf{หน่วยงาน} & \textbf{โครงการ} & \textbf{งบจัดสรร} & \textbf{เบิกจ่าย} & \textbf{ร้อยละ}\\
\midrule
กลุ่มแผนงานใต้ร่มพระบารมี & 16 & 2,879,963 & 1,472,124 & 51\%\\
สถาบันถ่ายทอดเทคโนโลยีสู่ชุมชน & 13 & 1,423,650 & 635,692 & 45\%\\
มทร.ล้านนา เชียงราย & 7 & 1,118,960 & 430,405 & 38\%\\
คณะวิศวกรรมศาสตร์ & 6 & 688,000 & 341,714 & 50\%\\
มทร.ล้านนา น่าน & 4 & 555,450 & 514,250 & 93\%\\
วิทยาลัยเทคโนโลยีและสหวิทยาการ & 2 & 349,060 & 255,349 & 73\%\\
มทร.ล้านนา พิษณุโลก & 5 & 331,500 & 233,928 & 71\%\\
มทร.ล้านนา ตาก & 3 & 250,000 & 0 & 0\%\\
คณะศิลปกรรมและสถาปัตยกรรมศาสตร์ & 2 & 200,000 & 92,144 & 46\%\\
สถาบันวิจัยเทคโนโลยีเกษตร & 2 & 150,000 & 42,198 & 28\%\\
คณะบริหารธุรกิจและศิลปศาสตร์ & 1 & 40,000 & 0 & 0\%\\
\midrule
\rowcolor{rpfgold!15}
\textbf{รวม} & \textbf{61} & \textbf{7,986,583} & \textbf{4,017,804} & \textbf{50\%}\\
\bottomrule
\end{tabular}
\end{center}


\keybox{rpfred}{\textbf{หน่วยงานที่ยังไม่มีการเบิกจ่าย} --- มทร.ล้านนา ตาก 3 โครงการ วงเงิน 250,000 บาท
และคณะบริหารธุรกิจและศิลปศาสตร์ 1 โครงการ วงเงิน 40,000 บาท ต้องเข้าไปสนับสนุนกระบวนการจัดซื้อจัดจ้างเร่งด่วน}

\clearpage
% ================= 4 ตัวชี้วัด =================
\section{ผลการดำเนินงานด้านตัวชี้วัด}
\subsection{ตัวชี้วัดหลักของมหาวิทยาลัยที่กลุ่มแผนงานรับมาดำเนินการ}

\begin{center}\small
\setlength{\tabcolsep}{5pt}\renewcommand{\arraystretch}{1.3}
\begin{tabular}{p{1.5cm}p{6.1cm}rrrr}
\toprule
\rowcolor{rpfnavy!12}
\textbf{ตัวชี้วัด} & \textbf{รายละเอียด} & \textbf{เป้า มทร.} & \textbf{รับมา} & \textbf{ร้อยละ} & \textbf{โครงการ}\\
\midrule
\rowcolor{rpfgreen!12}
\textbf{ที่ 39} & \textbf{องค์ความรู้/เทคโนโลยีที่นำไปใช้ในโครงการหลวงและโครงการตามพระราชดำริ (เป้าใต้ร่มโดยตรง)} & \textbf{60} & \textbf{62} & \textbf{103\%} & \textbf{37}\\
ที่ 35 & คณาจารย์และบุคลากรที่นำเทคโนโลยีไปพัฒนาชุมชนและภาคอุตสาหกรรม & 400 & 288 & 72\% & 46\\
ที่ 40 & องค์ความรู้และนวัตกรรมที่ยกระดับคุณภาพชีวิตชุมชน & 100 & 45 & 45\% & 33\\
ที่ 36 & แหล่งเรียนรู้ตลอดชีวิตของสังคม & 15 & 4 & 27\% & 4\\
\rowcolor{rpfred!8}
ที่ 10 & เครือข่ายความร่วมมือกับหน่วยงานภายนอก & 50 & 9 & 18\% & 4\\
\rowcolor{rpfred!8}
ที่ 17 & ทรัพย์สินทางปัญญาที่ยื่นขอจดทะเบียน & 60 & 8 & 13\% & 8\\
\rowcolor{rpfred!8}
ที่ 38 & สถานประกอบการที่ได้รับการถ่ายทอดเทคโนโลยี & 100 & 8 & 8\% & 7\\
\bottomrule
\end{tabular}
\end{center}

{\footnotesize \textbf{หมายเหตุ:} เป้าหมายในคอลัมน์ ``เป้า มทร.'' เป็นเป้าระดับมหาวิทยาลัย
กลุ่มแผนงานใต้ร่มพระบารมีเป็นผู้สนับสนุนส่วนหนึ่งเท่านั้น ยกเว้นตัวชี้วัดที่ 39 ซึ่งเป็นเป้าหมายของกลุ่มโดยตรง}

\subsection{การกระจายตัวชี้วัดตามโครงการหลัก ง8}

\begin{center}\small
\setlength{\tabcolsep}{6pt}\renewcommand{\arraystretch}{1.3}
\begin{tabular}{p{3.4cm}rrrrrr}
\toprule
\rowcolor{rpfnavy!12}
\textbf{โครงการหลัก} & \textbf{ที่ 10} & \textbf{ที่ 17} & \textbf{ที่ 35} & \textbf{ที่ 36} & \textbf{ที่ 38} & \textbf{ที่ 39 / 40}\\
\midrule
ง8-1 ผลักดันเทคโนโลยี & 1 & \textbf{8} & 98 & -- & 2 & 26 / 13\\
ง8-2 ขับเคลื่อนองค์ความรู้ & 1 & -- & 58 & \textbf{4} & 1 & 13 / 10\\
ง8-3 พัฒนากำลังคน & 7 & -- & \textbf{132} & -- & 5 & 23 / 22\\
\midrule
\rowcolor{rpfgold!15}
\textbf{รวม} & \textbf{9} & \textbf{8} & \textbf{288} & \textbf{4} & \textbf{8} & \textbf{62 / 45}\\
\bottomrule
\end{tabular}
\end{center}

ข้อสังเกต: ทรัพย์สินทางปัญญาทั้งหมด (ตัวชี้วัดที่ 17) อยู่ในโครงการหลัก ง8-1 ซึ่งเป็นแกนด้านเทคโนโลยี
แหล่งเรียนรู้ทั้งหมด (ตัวชี้วัดที่ 36) อยู่ใน ง8-2 และการพัฒนากำลังคน (ตัวชี้วัดที่ 35) กระจุกตัวใน ง8-3
สะท้อนว่าการแบ่งบทบาทของสามโครงการหลักมีความชัดเจน

\clearpage
% ================= 5 ผลราย ง8 =================
\section{ผลการดำเนินงานรายโครงการหลัก}

\subsection{ง8-1 ผลักดันเทคโนโลยี นวัตกรรมสู่ชุมชน ตามเป้าหมายการพัฒนาอย่างยั่งยืน}
16 โครงการย่อย $\cdot$ 6 หน่วยงานร่วม $\cdot$ งบจัดสรร 1,988,000 บาท $\cdot$ เบิกจ่าย 943,416 บาท (48\%)

\begin{center}\small
\setlength{\tabcolsep}{6pt}\renewcommand{\arraystretch}{1.3}
\begin{tabular}{p{2.5cm}p{12.2cm}}
\toprule
\rowcolor{rpfnavy!12}\textbf{ระดับผล} & \textbf{สาระสำคัญ}\\
\midrule
\textbf{ผลผลิต} & เครื่องต้นแบบและเทคโนโลยีการผลิตและแปรรูปบนพื้นที่สูง ; ทรัพย์สินทางปัญญาที่จะยื่นจดทะเบียน 8 ผลงาน ; องค์ความรู้ส่งมอบพื้นที่โครงการหลวง 26 รายการ\\
\textbf{ผลลัพธ์} & ศูนย์พัฒนาโครงการหลวงและชุมชนใช้เทคโนโลยีลดต้นทุนและเพิ่มคุณภาพผลผลิต ; บุคลากร 98 คน นำเทคโนโลยีลงพื้นที่จริง\\
\textbf{ผลกระทบ} & ยกระดับรายได้และขีดความสามารถของเกษตรกรบนพื้นที่สูง ; สร้างฐานทรัพย์สินทางปัญญาให้มหาวิทยาลัย\\
\bottomrule
\end{tabular}
\end{center}

\subsection{ง8-2 ขับเคลื่อนกลไกการพัฒนาองค์ความรู้เพื่อยกระดับคุณภาพชีวิต}
14 โครงการย่อย $\cdot$ 8 หน่วยงานร่วม $\cdot$ งบจัดสรร 1,999,010 บาท $\cdot$ เบิกจ่าย 1,276,760 บาท (66\% --- สูงสุดในแผนงาน)

\begin{center}\small
\setlength{\tabcolsep}{6pt}\renewcommand{\arraystretch}{1.3}
\begin{tabular}{p{2.5cm}p{12.2cm}}
\toprule
\rowcolor{rpfnavy!12}\textbf{ระดับผล} & \textbf{สาระสำคัญ}\\
\midrule
\textbf{ผลผลิต} & แหล่งเรียนรู้ตลอดชีวิต 4 แหล่ง ; พื้นที่เรียนรู้ศาสตร์พระราชาและเกษตรทฤษฎีใหม่ ; องค์ความรู้ส่งมอบพื้นที่ 13 รายการ\\
\textbf{ผลลัพธ์} & ชุมชนเข้าถึงองค์ความรู้ผ่านแหล่งเรียนรู้ถาวร ; บุคลากร 58 คน ถ่ายทอดความรู้ในพื้นที่ ; ศูนย์พัฒนาพันธุ์พืชจักรพันธ์เพ็ญศิริและงานรับเสด็จได้รับการสนับสนุนต่อเนื่อง\\
\textbf{ผลกระทบ} & เกิดกลไกการเรียนรู้ที่ยั่งยืนไม่ผูกกับรอบปีงบประมาณ ; คุณภาพชีวิตชุมชนดีขึ้นตามหลักปรัชญาเศรษฐกิจพอเพียง\\
\bottomrule
\end{tabular}
\end{center}

\subsection{ง8-3 พัฒนากำลังคน สร้างอาชีพ ลดความเหลื่อมล้ำ บนพื้นที่สูง}
31 โครงการย่อย $\cdot$ 9 หน่วยงานร่วม $\cdot$ งบจัดสรร 3,999,573 บาท (50\% ของงบแผนงาน) $\cdot$ เบิกจ่าย 1,797,628 บาท (46\%)

\begin{center}\small
\setlength{\tabcolsep}{6pt}\renewcommand{\arraystretch}{1.3}
\begin{tabular}{p{2.5cm}p{12.2cm}}
\toprule
\rowcolor{rpfnavy!12}\textbf{ระดับผล} & \textbf{สาระสำคัญ}\\
\midrule
\textbf{ผลผลิต} & หลักสูตรและการอบรมทักษะอาชีพ Re-skill Up-skill New-skill ; เครือข่ายความร่วมมือ 7 เครือข่าย ; องค์ความรู้ส่งมอบพื้นที่ 23 รายการ\\
\textbf{ผลลัพธ์} & เยาวชน สามเณร เจ้าหน้าที่ และชาวบ้านมีทักษะอาชีพใหม่ ; บุคลากร 132 คน ลงพื้นที่พัฒนากำลังคน ; วิสาหกิจชุมชนยกระดับมาตรฐานผลิตภัณฑ์\\
\textbf{ผลกระทบ} & ลดความเหลื่อมล้ำทางโอกาสและรายได้บนพื้นที่สูง ; สร้างกำลังคนที่อยู่กับชุมชนได้ในระยะยาว\\
\bottomrule
\end{tabular}
\end{center}

\clearpage
% ================= 6 วิเคราะห์ =================
\section{วิเคราะห์ผลการดำเนินงาน}

\subsection{จุดแข็ง}
\begin{enumerate}[leftmargin=1.6em,itemsep=1pt]
  \item \textbf{ภารกิจหลักตามพระราชดำริทำได้เกินเป้า} --- ตัวชี้วัดที่ 39 อยู่ที่ 103\% (62 จาก 60) กระจายใน 37 โครงการ สะท้อนว่าโครงการตอบโจทย์พื้นที่จริง
  \item \textbf{การกระจายภารกิจครอบคลุมทั้งมหาวิทยาลัย} --- 11 หน่วยงาน ทั้งคณะ สถาบัน วิทยาลัย และวิทยาเขต เชียงราย น่าน ตาก พิษณุโลก
  \item \textbf{กลไกอนุมัติและโอนงบทำงานได้ดี} --- อนุมัติแล้ว 59 จาก 61 โครงการ โอนงบถึงผู้รับผิดชอบ 98\%
  \item \textbf{มีระบบติดตามกลางแล้ว} --- ฐานข้อมูลเชื่อมรหัส ERP กับโครงการ แผนงาน หน่วยงาน และตัวชี้วัด พร้อมแดชบอร์ดผู้บริหาร
\end{enumerate}

\subsection{ช่องว่างและความเสี่ยง}
\begin{enumerate}[leftmargin=1.6em,itemsep=1pt]
  \item \textbf{ความเสี่ยงด้านการเบิกจ่าย (สูง)} --- เหลือเวลาประมาณ 70 วัน ต้องเบิกจ่ายอีก 3,968,779 บาท
        โดยมีภาระผูกพันรองรับแล้วเพียง 730,948 บาท หมายความว่าอีกกว่า 3.2 ล้านบาทยังไม่เริ่มกระบวนการจัดซื้อจัดจ้าง
  \item \textbf{ความเสี่ยงด้านข้อมูล (สูง)} --- ยังไม่มีรายงานผลรายกิจกรรมเข้าระบบ (0 จาก 268 กิจกรรม)
        ทำให้ไม่สามารถยืนยันผลสำเร็จที่ตรวจรับได้ และรายงานตัวชี้วัดต้องอ้างค่าที่รับมาดำเนินการแทน
  \item \textbf{ช่องว่างเชิงออกแบบตัวชี้วัด (ปานกลาง)} --- ตัวชี้วัดที่ 10 17 และ 38 ต่ำกว่า 20\%
        เนื่องจากโครงการส่วนใหญ่ออกแบบเพื่อชุมชนและพื้นที่โครงการหลวง มิได้ออกแบบเพื่อสถานประกอบการ
        หรือเพื่อการจดทะเบียนทรัพย์สินทางปัญญาตั้งแต่ต้น
  \item \textbf{ความเหลื่อมของอัตราเบิกจ่ายระหว่างหน่วยงาน (ปานกลาง)} --- ตั้งแต่ 93\% (มทร.ล้านนา น่าน)
        ลงมาถึง 0\% (มทร.ล้านนา ตาก และคณะบริหารธุรกิจและศิลปศาสตร์)
\end{enumerate}

\keybox{rpfnavy}{\textbf{ข้อสรุปเชิงวิเคราะห์} --- ปัญหาที่พบเป็นเรื่อง\textbf{กระบวนการเบิกจ่าย การรายงานข้อมูล
และการออกแบบตัวชี้วัด} มิใช่เรื่องคุณภาพของงานในพื้นที่ ซึ่งสะท้อนได้จากตัวชี้วัดภารกิจหลักที่ทำได้เกินเป้าแล้ว}

\clearpage
% ================= 7 แผนต่อ =================
\section{แผนการดำเนินงานในช่วงที่เหลือของปีงบประมาณ}

\begin{center}\small
\setlength{\tabcolsep}{6pt}\renewcommand{\arraystretch}{1.35}
\begin{tabular}{p{1.9cm}p{9.4cm}p{3.4cm}}
\toprule
\rowcolor{rpfnavy!12}
\textbf{ช่วงเวลา} & \textbf{กิจกรรม} & \textbf{เป้าหมาย}\\
\midrule
สิงหาคม 2569 & เร่งรัดการเบิกจ่ายและปิดใบสั่งซื้อคงค้าง ; ลงพื้นที่ติดตามผลและเก็บหลักฐานตัวชี้วัด ; สนับสนุนหน่วยงานที่ยังไม่เบิกจ่าย & เบิกจ่ายสะสมถึง 75\%\\
กันยายน 2569 (ต้น--กลาง) & บันทึกรายงานผลรายกิจกรรมเข้าระบบให้ครบทุกโครงการ ; ยื่นจดทะเบียนทรัพย์สินทางปัญญา ; ลงนามบันทึกความร่วมมือกับหน่วยงานภายนอก & รายงานผลเข้าระบบครบ 268 กิจกรรม\\
กันยายน 2569 (ปลาย) & ตรวจรับผลตัวชี้วัดจากหลักฐานจริง ; จัดทำรายงานผลการดำเนินงานฉบับสมบูรณ์ ; ถอดบทเรียนและจัดทำข้อเสนอปีงบประมาณ 2570 & ปิดปีงบประมาณ 30 กันยายน 2569\\
\bottomrule
\end{tabular}
\end{center}

\subsection{มาตรการเร่งรัด 4 ด้าน}
\begin{enumerate}[leftmargin=1.6em,itemsep=2pt]
  \item \textbf{เร่งรัดการเบิกจ่าย} --- ติดตามรายโครงการที่เบิกต่ำกว่า 30\% เป็นรายสัปดาห์ และช่วยจัดทำเอกสารจัดซื้อจัดจ้างให้หน่วยงานที่ยังไม่เบิกจ่าย
  \item \textbf{เก็บผลงานเข้าระบบ} --- กำหนดให้ผู้รับผิดชอบรายงานผลรายกิจกรรมพร้อมหลักฐานภาพถ่ายและเอกสารในระบบติดตาม แล้วตรวจรับยืนยันผลตัวชี้วัดจากหลักฐานจริง
  \item \textbf{ปิดช่องว่างตัวชี้วัด} --- เร่งยื่นจดทรัพย์สินทางปัญญาจากผลงานที่พร้อม และจัดทำบันทึกความร่วมมือกับหน่วยงานและสถานประกอบการในพื้นที่
  \item \textbf{เตรียมต่อยอดปี 2570} --- ถอดบทเรียนรายพื้นที่ และออกแบบตัวชี้วัดให้สอดคล้องกับลักษณะงานพัฒนาชุมชนตั้งแต่ต้นปีงบประมาณ
\end{enumerate}

\subsection{การตอบโจทย์เป้าหมายการพัฒนาที่ยั่งยืน}
ในภาพรวมระดับแผนงาน การดำเนินงานเชื่อมโยงกับเป้าหมายการพัฒนาที่ยั่งยืน (SDGs) หลัก ได้แก่
SDG 1 ขจัดความยากจน $\cdot$ SDG 2 ยุติความหิวโหย $\cdot$ SDG 4 การศึกษาที่มีคุณภาพ $\cdot$
SDG 8 งานที่มีคุณค่าและการเติบโตทางเศรษฐกิจ $\cdot$ SDG 9 อุตสาหกรรม นวัตกรรม และโครงสร้างพื้นฐาน $\cdot$
SDG 10 ลดความเหลื่อมล้ำ $\cdot$ SDG 13 การรับมือการเปลี่ยนแปลงสภาพภูมิอากาศ $\cdot$
SDG 17 ความร่วมมือเพื่อการพัฒนาที่ยั่งยืน
และเป้าหมายสนับสนุน ได้แก่ SDG 7 พลังงานสะอาด $\cdot$ SDG 11 ชุมชนยั่งยืน $\cdot$
SDG 12 การผลิตและบริโภคที่ยั่งยืน $\cdot$ SDG 15 ระบบนิเวศบนบก

{\footnotesize \textbf{หมายเหตุ:} การเชื่อมโยงเป้าหมายการพัฒนาที่ยั่งยืนข้างต้นเป็นการวิเคราะห์ระดับแผนงาน
จากวัตถุประสงค์ของโครงการหลัก ยังมิได้ผูกรายโครงการไว้ในฐานข้อมูล}

\clearpage
% ================= ภาคผนวก ก =================
\section*{ภาคผนวก ก --- รายชื่อโครงการและการเบิกจ่ายรายโครงการ}
\addcontentsline{toc}{section}{ภาคผนวก ก --- รายชื่อโครงการและการเบิกจ่ายรายโครงการ}
\markboth{}{}

{\footnotesize
\setlength{\tabcolsep}{3.5pt}\renewcommand{\arraystretch}{1.22}
\begin{longtable}{@{}p{0.5cm}p{6.1cm}p{2.9cm}p{2.7cm}rr@{}}
\toprule
\rowcolor{rpfnavy!12}
\textbf{ที่} & \textbf{ชื่อโครงการ} & \textbf{ผู้รับผิดชอบ} & \textbf{หน่วยงาน} & \textbf{จัดสรร} & \textbf{เบิกจ่าย}\\
\midrule
\endfirsthead
\toprule
\rowcolor{rpfnavy!12}
\textbf{ที่} & \textbf{ชื่อโครงการ} & \textbf{ผู้รับผิดชอบ} & \textbf{หน่วยงาน} & \textbf{จัดสรร} & \textbf{เบิกจ่าย}\\
\midrule
\endhead
\multicolumn{6}{@{}l}{\cellcolor{rpfnavy!12}\bfseries ง8-1 ผลักดันเทคโนโลยี นวัตกรรมสู่ชุมชน ตามเป้าหมายการพัฒนาอย่างยั่งยืน (16 โครงการ)}\\
\midrule
1 & โครงการการวิจัย Plant profile ของผักในโรงเรือนอัจฉริยะ & นายวัชระ กิตติวรเชฏฐ์ & กลุ่มแผนงานใต้ร่มพระบารมี & 200,000 & 46,387 (23\%)\\
2 & โครงการการวิจัยพัฒนาวิธีการจัดการกระบวนการผลิตผักเพื่อคุณภาพ & นายวัชระ กิตติวรเชฏฐ์ & กลุ่มแผนงานใต้ร่มพระบารมี & 150,000 & 61,222 (41\%)\\
3 & โครงการพัฒนากระบวนการสกัดน้ำมันอะโวคาโดของศูนย์พัฒนาโครงการหลวงห้วยเสี & นางสาวพิมลพรรณ เลิศบัวบา & กลุ่มแผนงานใต้ร่มพระบารมี & 148,500 & 63,002 (42\%)\\
4 & โครงการยกระดับการพัฒนาเทคโนโลยี เครื่องมือ และเครื่องจักรของมหาวิทยาลั & นายสามารถ สาลี & กลุ่มแผนงานใต้ร่มพระบารมี & 52,000 & 0 (0\%)\\
5 & โครงการการพัฒนาตู้อบลมร้อนดอกเก๊กฮวยระบบปั๊มความร้อน เพื่อเพิ่มคุณภาพผ & ผศ.ไกรลาศ ดอนชัย & คณะวิศวกรรมศาสตร์ & 175,000 & 87,500 (50\%)\\
6 & โครงการพัฒนาเทคโนโลยีเครื่องล้างเมล็ดงาดำแบบกึ่งอัตโนมัติ สำหรับศูนย์พ & ผศ.ขัติพงษ์ จิโนสุวัตร์ & คณะวิศวกรรมศาสตร์ & 150,000 & 67,432 (45\%)\\
7 & โครงการพัฒนาเทคโนโลยีต้นแบบเครื่องนวดงาขี้ม่อน(งาหอม) เพื่อเพิ่มประสิท & อ.วรเชษฐ์ หวานเสียง & คณะวิศวกรรมศาสตร์ & 175,000 & 160,492 (92\%)\\
8 & โครงการประเมินสมบัติเชิงกลของเส้นใยกัญชงจากการอบแห้งด้วยเทคโนโลยีอบลมร & นายธีระยุทธ ขอดแก้ว & คณะวิศวกรรมศาสตร์ & 50,000 & 23,200 (46\%)\\
9 & โครงการศึกษาความเหมาะสมของการนำเศษผักเหลือใช้มาผลิตเป็นถ่านอัดแท่งขนาด & นายเมธัส ภัททิยธนี & คณะวิศวกรรมศาสตร์ & 88,000 & 0 (0\%)\\
10 & โครงการ Royal Rose Sensation Giftset ชุดผลิตภัณฑ์ต่อยอดคุณค่ากุหลาบจาก & อ.ณัฐธินี สาลี & วิทยาลัยเทคโนโลยีและสหวิทยาการ & 200,000 & 185,349 (93\%)\\
11 & โครงการพัฒนากระบวนการผลิตและเครื่องมือการอบแห้งเคพกูสเบอร์รี เพื่อเพิ่ & นายวธัญญู วรรณพรหม & สถาบันถ่ายทอดเทคโนโลยีสู่ชุมชน & 100,000 & 20,000 (20\%)\\
12 & โครงการพัฒนาหลักสูตรชุมชนการเรียนรู้เพื่อการพัฒนาตนเองตามภูมิสั่งคมอย่ & ผู้ช่วยศาสตราจารย์ไพโรจน & สถาบันถ่ายทอดเทคโนโลยีสู่ชุมชน & 68,000 & 32,200 (47\%)\\
13 & โครงการพัฒนาชุดทดสอบสารพาราควอตเบื้องต้น (Paraquat Test Kit) ที่เหมาะส & ผศ.สุบิน ใจทา & มทร.ล้านนา เชียงราย & 150,000 & 20,000 (13\%)\\
14 & โครงการนวัตกรรมเชิงพื้นที่ของกระถางเพาะเมล็ดชีวภาพย่อยสลายได้จากของเหล & ผศ.อมรรัตน์ ปิ่นชัยมูล & มทร.ล้านนา เชียงราย & 100,000 & 60,268 (60\%)\\
15 & โครงการประยุกต์ใช้เทคโนโลยีที่เหมาะสมเพื่อเพิ่มประสิทธิภาพการเลี้ยงปลา & นางสาวอพิศรา หงส์หิรัญ & มทร.ล้านนา พิษณุโลก & 100,000 & 76,584 (77\%)\\
16 & โครงการพัฒนาเครื่องสลัดน้ำผักตามอัตลักษณ์ของศูนย์พัฒนาโครงการหลวงทุ่งห & ผศ.เดือนแรม แพ่งเกี่ยว & มทร.ล้านนา พิษณุโลก & 81,500 & 39,780 (49\%)\\
\midrule
\multicolumn{6}{@{}l}{\cellcolor{rpfnavy!12}\bfseries ง8-2 ขับเคลื่อนกลไกการพัฒนาองค์ความรู้เพื่อยกระดับคุณภาพชีวิต (14 โครงการ)}\\
\midrule
1 & โครงการพัฒนาพื้นที่เรียนรู้ศาสตร์พระราชา เพื่อยกระดับศักยภาพพื้นที่ด้ว & นายสามารถ สาลี & กลุ่มแผนงานใต้ร่มพระบารมี & 400,000 & 285,850 (71\%)\\
2 & โครงการการถ่ายทอดองค์ความรู้เรื่องการจัดการขยะแบบยั่งยืน ระดับชุมชนและ & ดร.วนิดา สุริยานนท์ & คณะวิศวกรรมศาสตร์ & 50,000 & 3,090 (6\%)\\
3 & โครงการอบรมเชิงปฏิบัติการถ่ายทอดองค์ความรู้เทคนิคการสร้างแม่พิมพ์เพื่อ & นายคเชนทร์ เครือสาร & สถาบันถ่ายทอดเทคโนโลยีสู่ชุมชน & 50,000 & 47,751 (96\%)\\
4 & โครงการการพัฒนาแหล่งเรียนรู้และถ่ายทอดองค์ความรู้ เทคโนโลยีการพัฒนาผลิ & นายศรีธร อุปคำ & สถาบันถ่ายทอดเทคโนโลยีสู่ชุมชน & 200,000 & 111,400 (56\%)\\
5 & โครงการบริหารจัดการโครงการผลิตเมล็ดพันธุ์และปรับปรุงพันธุ์พืชให้กับศูน & นายศิริชัย แซ่ท้าว & สถาบันวิจัยเทคโนโลยีเกษตร & 100,000 & 42,198 (42\%)\\
6 & โครงการการจัดนิทรรศการเพื่อน้อมถวายรายงานในการรับเสด็จสมเด็จพระกนิษฐาธ & นางสาวสุภัควดี พิมพ์มาศ & มทร.ล้านนา น่าน & 250,000 & 250,000 (100\%)\\
7 & โครงการพัฒนาศักยภาพการเรียนรู้ ด้านการจัดการภูมิทัศน์ ผ่านประสบการณ์จร & นายภาณุพงศ์ สิทธิวุฒิ & มทร.ล้านนา น่าน & 129,050 & 129,050 (100\%)\\
8 & โครงการสนับสนุนการดำเนินงานศูนย์พัฒนาพันธุ์พืชจักรพันธ์เพ็ญศิริ อ.แม่ส & รศ.วิเชษฐ ทิพย์ประเสริฐ & มทร.ล้านนา เชียงราย & 450,000 & 279,847 (62\%)\\
9 & โครงการพัฒนาแหล่งเรียนรู้เกษตรอันเนื่องมาจากพระราชดำริ วิถีธรรมชาติ & รศ.วิเชษฐ ทิพย์ประเสริฐ & มทร.ล้านนา เชียงราย & 118,960 & 10,010 (8\%)\\
10 & โครงการจัดทำคำขอสิ่งบ่งชี้ทางภูมิศาสตร์ไทย (GI) สินค้ากาแฟเลอตอ (อาราบ & ผู้ช่วยศาสตราจารย์ทนงศัก & มทร.ล้านนา ตาก & 36,000 & 0 (0\%)\\
11 & โครงการพัฒนาพื้นที่ต้นแบบ โคก หนอง นา เพื่อการเรียนรู้ด้านเกษตรอินทรีย & ผู้ช่วยศาสตราจารย์ทนงศัก & มทร.ล้านนา ตาก & 65,000 & 0 (0\%)\\
12 & โครงการประยุกต์ใช้การคำนวนสูตรอาหารไก่ไข่ต้นทุนต่ำด้วยโปรแกรมสำเร็จรูป & นางสาวอุษณีย์ภรณ์ สร้อยเ & มทร.ล้านนา พิษณุโลก & 50,000 & 42,500 (85\%)\\
13 & โครงการถ่ายทอดองค์ความรู้และนวัตกรรมการยกระดับวัสดุเพาะเห็ดเศรษฐกิจจาก & นายภควัต หุ่นฉัตร์ & มทร.ล้านนา พิษณุโลก & 50,000 & 27,802 (56\%)\\
14 & โครงการเพิ่มมูลค่าสินค้าเกษตรและการพัฒนาสื่อด้วย AI พื้นที่เขตปฏิรูปที & ผศ.สุริยาพร นิพรรัมย์ & มทร.ล้านนา พิษณุโลก & 50,000 & 47,262 (95\%)\\
\midrule
\multicolumn{6}{@{}l}{\cellcolor{rpfnavy!12}\bfseries ง8-3 พัฒนากำลังคน สร้างอาชีพ ลดความเหลื่อมล้ำ บนพื้นที่สูง (31 โครงการ)}\\
\midrule
1 & โครงการพัฒนาการบริหารจัดการฐานข้อมูลด้วยกลไกการติดตามและประเมินผล โครง & พิมลพรรณ เลิศบัวบาน & กลุ่มแผนงานใต้ร่มพระบารมี & 500,000 & 349,316 (70\%)\\
2 & โครงการพัฒนาฐานข้อมูลด้านวิศวกรรมและยกระดับทักษะอาชีพในพื้นที่โครงการห & นายสามารถ สาลี & กลุ่มแผนงานใต้ร่มพระบารมี & 200,000 & 81,520 (41\%)\\
3 & โครงการยกระดับการบริหารจัดการ ระบบติดตามและเครือข่าย ความร่วมมือกลุ่มแ & นายนริศ กำแพงแก้ว & กลุ่มแผนงานใต้ร่มพระบารมี & 500,000 & 203,582 (41\%)\\
4 & โครงการการจัดกิจกรรมการเรียนรู้ส่งเสริมงานวิจัยและบริการวิชาการ ด้านนว & นายวัชระ กิตติวรเชฏฐ์ & กลุ่มแผนงานใต้ร่มพระบารมี & 60,563 & 60,277 (100\%)\\
5 & โครงการการฝึกอบรมสร้างความสามารถการออกแบบระบบจ่ายน้ำในแปลงเกษตร & นายวัชระ กิตติวรเชฏฐ์ & กลุ่มแผนงานใต้ร่มพระบารมี & 150,000 & 117,038 (78\%)\\
6 & โครงการพัฒนาทักษะอาชีพด้านพลังงานสะอาด & นายวธัญญู วรรณพรหม & กลุ่มแผนงานใต้ร่มพระบารมี & 100,000 & 27,790 (28\%)\\
7 & โครงการนำร่อง SkillChain มทร.ล้านนา :ระบบรับรองทักษะช่างและจัดหางานบน & นางสาวพิมลพรรณ เลิศบัวบา & กลุ่มแผนงานใต้ร่มพระบารมี & 100,000 & 0 (0\%)\\
8 & โครงการการสำรวจผลผลิตทางการเกษตรโครงการหลวง เพื่อการประกอบอาหารและเผยแ & นายวรัญญ์คมน์ รัตนะมงคลช & คณะบริหารธุรกิจและศิลปศาสตร์ & 40,000 & 0 (0\%)\\
9 & โครงการยกระดับมาตรฐานผลิตภัณฑ์ผ้าใยกัญชงวิสาหกิจชุมชนดาวม่างสู่ มผช. ด & ผศ.ญาณิศา โกมลสิริโชค & คณะศิลปกรรมและสถาปัตยกรรมศาสตร์ & 100,000 & 49,644 (50\%)\\
10 & โครงการยกระดับห่วงโซ่คุณค่าห้อมพื้นถิ่น สู่ผลิตภัณฑ์ย้อมสี ธรรมชาติสร้ & ดร.จุฑามาศ ดอนอ่อนเบ้า & คณะศิลปกรรมและสถาปัตยกรรมศาสตร์ & 100,000 & 42,500 (43\%)\\
11 & โครงการการพัฒนานวัตกรรมด้านการศึกษาเพื่อการพัฒนาศักยภาพเยาวชนในโรงเรีย & นายศุภกาล ตุ้ยเต็มวงศ์ & วิทยาลัยเทคโนโลยีและสหวิทยาการ & 149,060 & 70,000 (47\%)\\
12 & โครงการประเมินผลกระทบของโครงการบริการวิชาการมหาวิทยาลัยเทคโนโลยีราชมงค & ผศ.ธัญวดี สุจริตธรรม & สถาบันถ่ายทอดเทคโนโลยีสู่ชุมชน & 142,800 & 90,890 (64\%)\\
13 & โครงการพัฒนาแผนการตลาดของกลุ่มวิสาหกิจชุมชนภายใต้แนวคิด “กินได้ ใช้ได้ & ผศ.ธัญวดี สุจริตธรรม & สถาบันถ่ายทอดเทคโนโลยีสู่ชุมชน & 63,000 & 33,197 (53\%)\\
14 & โครงการต้นแบบอาคารอนุรักษ์พลังงานในรูปแบบ (Green Building \& Green Offi & นายวธัญญู วรรณพรหม & สถาบันถ่ายทอดเทคโนโลยีสู่ชุมชน & 250,000 & 107,500 (43\%)\\
15 & โครงการพัฒนาทักษะวิชาชีพให้กับเจ้าหน้าที่ ชาวบ้าน เยาวชน พื้นที่โดยรอบ & นายศราวุธ วรรณวิจิตร & สถาบันถ่ายทอดเทคโนโลยีสู่ชุมชน & 100,000 & 31,875 (32\%)\\
16 & โครงการพัฒนาต้นแบบชุดน้ำชาเซรามิกสโตนแวร์สำหรับยกระดับผลิตภัณฑ์ชาโครงก & นายสิงหล วิชายะ & สถาบันถ่ายทอดเทคโนโลยีสู่ชุมชน & 100,000 & 0 (0\%)\\
17 & โครงการการถ่ายทอดองค์ความรู้กระบวนการผลิตน้ำดื่มสะอาดด้วยระบบรถน้ำดื่ม & นายเมธัส ภัททิยธนี & สถาบันถ่ายทอดเทคโนโลยีสู่ชุมชน & 150,000 & 10,483 (7\%)\\
18 & โครงการพัฒนาทักษะวิชาชีพให้กับเจ้าหน้าที่ ชาวบ้าน เยาวชน พื้นที่โดยรอบ & นายวธัญญู วรรณพรหม & สถาบันถ่ายทอดเทคโนโลยีสู่ชุมชน & 100,000 & 73,386 (73\%)\\
19 & การประยุกต์ใช้ lora mesh ชนิด multi-hops ในการส่งสัญญาณจากเซนเซอร์ตรวจ & นายณัฐวัฒน์ พัลวัล & กลุ่มแผนงานใต้ร่มพระบารมี & 100,000 & 72,800 (73\%)\\
20 & การประยุกต์ใช้เครือข่าย lora-mesh ชนิด multi-hops เพื่อการส่งสัญญาณเตื & นายณัฐวัฒน์ พัลวัล & กลุ่มแผนงานใต้ร่มพระบารมี & 100,000 & 57,500 (57\%)\\
21 & โครงการพัฒนาและยกระดับสมรรถนะด้านการติดตั้งและบำรุงรักษาระบบเซลล์แสงอา & นายวธัญญู วรรณพรหม & กลุ่มแผนงานใต้ร่มพระบารมี & 18,900 & 18,900 (100\%)\\
22 & โครงการวิจัยความเข้มแข็งทางสังคม ชุมชนในการวางแผนการใช้น้ำศูนย์พัฒนาโค & นายวัชระ กิตติวรเชฏฐ์ & กลุ่มแผนงานใต้ร่มพระบารมี & 100,000 & 26,940 (27\%)\\
23 & โครงการส่งเสริมบำรุงรักษาระบบผลิตไฟฟ้าแบบผสมผสานระหว่างระบบผลิตไฟฟ้าพล & นายศรีธร อุปคำ & สถาบันถ่ายทอดเทคโนโลยีสู่ชุมชน & 49,990 & 48,150 (96\%)\\
24 & โครงการพัฒนาทักษะวิชาชีพหลักสูตรการทำมุ้งลวดให้กับสามเณร และเยาวชน พื้ & นายวธัญญู วรรณพรหม & สถาบันถ่ายทอดเทคโนโลยีสู่ชุมชน & 49,860 & 28,860 (58\%)\\
25 & โครงการขับเคลื่อนการสร้างกลไกการส่งเสริมและพัฒนาศักยภาพการเรียนรู้ด้าน & นายรัตนพล พนมวัน ณ อยุธย & สถาบันวิจัยเทคโนโลยีเกษตร & 50,000 & 0 (0\%)\\
26 & โครงการพัฒนาทักษะอาชีพและองค์ความรู้ด้านเกษตรและเทคโนโลยีอย่างยั่งยืน & นายก้องเกียรติ ธนะมิตร & มทร.ล้านนา น่าน & 94,000 & 94,000 (100\%)\\
27 & โครงการถ่ายทอดเทคโนโลยีการเพาะขยายพันธุ์ปลาเลียหินเพื่อการอนุรักษ์ปลาท & นายจุลทรรศน์ คีรีแลง & มทร.ล้านนา น่าน & 82,400 & 41,200 (50\%)\\
28 & โครงส่งเสริมพัฒนาโจทย์วิจัยจากชุมชน & รศ.วิเชษฐ ทิพย์ประเสริฐ & มทร.ล้านนา เชียงราย & 100,000 & 8,240 (8\%)\\
29 & โครงการถ่ายทอดองค์ความรู้และนวัตกรรมเพื่อยกระดับมูลค่าไก่ดำ (ไก่ฟ้าหลว & ผศ.ดร.จิรพัฒน์พงษ์ เสนาบ & มทร.ล้านนา เชียงราย & 100,000 & 7,040 (7\%)\\
30 & โครงการพัฒนาระบบติดตามข้อมูลภายใต้สถานีตรวจวัดสภาพอากาศเพื่อป้องกันภัย & ผศ.กิ่งกาญจน์ ปวนสุรินทร & มทร.ล้านนา เชียงราย & 100,000 & 45,000 (45\%)\\
31 & โครงการจัดทำคำขอสิ่งบ่งชี้ทางภูมิศาสตร์ไทย (GI) สินค้ากาแฟเลอตอ (อาราบ & ผู้ช่วยศาสตราจารย์ทนงศัก & มทร.ล้านนา ตาก & 149,000 & 0 (0\%)\\
\midrule
\bottomrule
\end{longtable}
}


\clearpage
% ================= ภาคผนวก ข =================
\section*{ภาคผนวก ข --- ตัวชี้วัดที่รับมาดำเนินการรายโครงการ}
\addcontentsline{toc}{section}{ภาคผนวก ข --- ตัวชี้วัดที่รับมาดำเนินการรายโครงการ}

{\footnotesize รูปแบบการอ่าน: \textbf{หมายเลขตัวชี้วัด:จำนวนที่รับมา} เช่น \textbf{39:2} หมายถึง ตัวชี้วัดที่ 39 จำนวน 2 องค์ความรู้\par\vspace{4pt}}

{\footnotesize
\setlength{\tabcolsep}{4pt}\renewcommand{\arraystretch}{1.22}
\begin{longtable}{@{}p{0.5cm}p{9.1cm}p{5.2cm}@{}}
\toprule
\rowcolor{rpfnavy!12}
\textbf{ที่} & \textbf{ชื่อโครงการ} & \textbf{ตัวชี้วัดที่รับมา}\\
\midrule
\endfirsthead
\toprule
\rowcolor{rpfnavy!12}
\textbf{ที่} & \textbf{ชื่อโครงการ} & \textbf{ตัวชี้วัดที่รับมา}\\
\midrule
\endhead
\multicolumn{3}{@{}l}{\cellcolor{rpfnavy!12}\bfseries ง8-1 ผลักดันเทคโนโลยี นวัตกรรมสู่ชุมชน ตามเป้าหมายการพัฒนาอย่างยั่งยืน}\\
\midrule
1 & โครงการการวิจัย Plant profile ของผักในโรงเรือนอัจฉริยะ & 10:1 $\cdot$ 35:2 $\cdot$ 39:1\\
2 & โครงการการวิจัยพัฒนาวิธีการจัดการกระบวนการผลิตผักเพื่อคุณภาพ & 17:1 $\cdot$ 35:2 $\cdot$ 39:2 $\cdot$ 40:1\\
3 & โครงการพัฒนากระบวนการสกัดน้ำมันอะโวคาโดของศูนย์พัฒนาโครงการหลวงห & 35:7 $\cdot$ 39:2 $\cdot$ 40:1\\
4 & โครงการการพัฒนาตู้อบลมร้อนดอกเก๊กฮวยระบบปั๊มความร้อน เพื่อเพิ่มค & 17:1 $\cdot$ 35:8 $\cdot$ 39:2 $\cdot$ 40:1\\
5 & โครงการพัฒนาเทคโนโลยีเครื่องล้างเมล็ดงาดำแบบกึ่งอัตโนมัติ สำหรับ & 17:1 $\cdot$ 35:12 $\cdot$ 39:2 $\cdot$ 40:1\\
6 & โครงการพัฒนาเทคโนโลยีต้นแบบเครื่องนวดงาขี้ม่อน(งาหอม) เพื่อเพิ่ม & 17:1 $\cdot$ 35:12 $\cdot$ 39:2 $\cdot$ 40:1\\
7 & โครงการประเมินสมบัติเชิงกลของเส้นใยกัญชงจากการอบแห้งด้วยเทคโนโลย & 35:8\\
8 & โครงการศึกษาความเหมาะสมของการนำเศษผักเหลือใช้มาผลิตเป็นถ่านอัดแท & 35:5 $\cdot$ 39:3 $\cdot$ 40:2\\
9 & โครงการ Royal Rose Sensation Giftset ชุดผลิตภัณฑ์ต่อยอดคุณค่ากุห & 17:1 $\cdot$ 35:15 $\cdot$ 38:1 $\cdot$ 39:2 $\cdot$ 40:1\\
10 & โครงการพัฒนากระบวนการผลิตและเครื่องมือการอบแห้งเคพกูสเบอร์รี เพื & 17:1 $\cdot$ 35:3 $\cdot$ 39:2 $\cdot$ 40:1\\
11 & โครงการพัฒนาชุดทดสอบสารพาราควอตเบื้องต้น (Paraquat Test Kit) ที่ & 17:1 $\cdot$ 35:4 $\cdot$ 39:1\\
12 & โครงการนวัตกรรมเชิงพื้นที่ของกระถางเพาะเมล็ดชีวภาพย่อยสลายได้จาก & 35:6 $\cdot$ 39:2 $\cdot$ 40:1\\
13 & โครงการประยุกต์ใช้เทคโนโลยีที่เหมาะสมเพื่อเพิ่มประสิทธิภาพการเลี & 35:6 $\cdot$ 38:1 $\cdot$ 39:3 $\cdot$ 40:2\\
14 & โครงการพัฒนาเครื่องสลัดน้ำผักตามอัตลักษณ์ของศูนย์พัฒนาโครงการหลว & 17:1 $\cdot$ 35:8 $\cdot$ 39:2 $\cdot$ 40:1\\
\midrule
\multicolumn{3}{@{}l}{\cellcolor{rpfnavy!12}\bfseries ง8-2 ขับเคลื่อนกลไกการพัฒนาองค์ความรู้เพื่อยกระดับคุณภาพชีวิต}\\
\midrule
1 & โครงการการถ่ายทอดองค์ความรู้เรื่องการจัดการขยะแบบยั่งยืน ระดับชุ & 35:4\\
2 & โครงการอบรมเชิงปฏิบัติการถ่ายทอดองค์ความรู้เทคนิคการสร้างแม่พิมพ & 35:6 $\cdot$ 39:1 $\cdot$ 40:1\\
3 & โครงการบริหารจัดการโครงการผลิตเมล็ดพันธุ์และปรับปรุงพันธุ์พืชให้ & 35:8 $\cdot$ 36:1\\
4 & โครงการพัฒนาศักยภาพการเรียนรู้ ด้านการจัดการภูมิทัศน์ ผ่านประสบก & 35:7\\
5 & โครงการสนับสนุนการดำเนินงานศูนย์พัฒนาพันธุ์พืชจักรพันธ์เพ็ญศิริ & 35:5 $\cdot$ 36:1 $\cdot$ 39:3 $\cdot$ 40:3\\
6 & โครงการพัฒนาแหล่งเรียนรู้เกษตรอันเนื่องมาจากพระราชดำริ วิถีธรรมช & 35:7 $\cdot$ 35:5 $\cdot$ 36:1 $\cdot$ 39:3 $\cdot$ 39:1 $\cdot$ 40:1 $\cdot$ 40:1\\
7 & โครงการจัดทำคำขอสิ่งบ่งชี้ทางภูมิศาสตร์ไทย (GI) สินค้ากาแฟเลอตอ & 38:1 $\cdot$ 39:1 $\cdot$ 40:1\\
8 & โครงการพัฒนาพื้นที่ต้นแบบ โคก หนอง นา เพื่อการเรียนรู้ด้านเกษตรอ & 10:1 $\cdot$ 36:1 $\cdot$ 39:1\\
9 & โครงการประยุกต์ใช้การคำนวนสูตรอาหารไก่ไข่ต้นทุนต่ำด้วยโปรแกรมสำเ & 35:4 $\cdot$ 39:1 $\cdot$ 40:1\\
10 & โครงการถ่ายทอดองค์ความรู้และนวัตกรรมการยกระดับวัสดุเพาะเห็ดเศรษฐ & 35:8 $\cdot$ 39:1 $\cdot$ 40:1\\
11 & โครงการเพิ่มมูลค่าสินค้าเกษตรและการพัฒนาสื่อด้วย AI พื้นที่เขตปฏ & 35:4 $\cdot$ 39:1 $\cdot$ 40:1\\
\midrule
\multicolumn{3}{@{}l}{\cellcolor{rpfnavy!12}\bfseries ง8-3 พัฒนากำลังคน สร้างอาชีพ ลดความเหลื่อมล้ำ บนพื้นที่สูง}\\
\midrule
1 & โครงการพัฒนาฐานข้อมูลด้านวิศวกรรมและยกระดับทักษะอาชีพในพื้นที่โค & 35:1 $\cdot$ 39:1 $\cdot$ 40:1\\
2 & โครงการยกระดับการบริหารจัดการ ระบบติดตามและเครือข่าย ความร่วมมือ & 10:5\\
3 & โครงการการจัดกิจกรรมการเรียนรู้ส่งเสริมงานวิจัยและบริการวิชาการ & 35:4\\
4 & โครงการการฝึกอบรมสร้างความสามารถการออกแบบระบบจ่ายน้ำในแปลงเกษตร & 35:1 $\cdot$ 39:1 $\cdot$ 40:2\\
5 & โครงการพัฒนาทักษะอาชีพด้านพลังงานสะอาด & 35:3 $\cdot$ 39:1 $\cdot$ 40:1\\
6 & โครงการนำร่อง SkillChain มทร.ล้านนา :ระบบรับรองทักษะช่างและจัดหา & 10:2\\
7 & โครงการการสำรวจผลผลิตทางการเกษตรโครงการหลวง เพื่อการประกอบอาหารแ & 35:8\\
8 & โครงการยกระดับมาตรฐานผลิตภัณฑ์ผ้าใยกัญชงวิสาหกิจชุมชนดาวม่างสู่ & 35:5 $\cdot$ 38:1 $\cdot$ 39:2 $\cdot$ 40:2\\
9 & โครงการยกระดับห่วงโซ่คุณค่าห้อมพื้นถิ่น สู่ผลิตภัณฑ์ย้อมสี ธรรมช & 35:9 $\cdot$ 38:1 $\cdot$ 39:2 $\cdot$ 40:1\\
10 & โครงการการพัฒนานวัตกรรมด้านการศึกษาเพื่อการพัฒนาศักยภาพเยาวชนในโ & 35:6 $\cdot$ 39:1 $\cdot$ 40:1\\
11 & โครงการประเมินผลกระทบของโครงการบริการวิชาการมหาวิทยาลัยเทคโนโลยี & 35:7\\
12 & โครงการพัฒนาแผนการตลาดของกลุ่มวิสาหกิจชุมชนภายใต้แนวคิด “กินได้ & 35:7 $\cdot$ 38:2 $\cdot$ 39:1\\
13 & โครงการต้นแบบอาคารอนุรักษ์พลังงานในรูปแบบ (Green Building \& Gree & 35:3\\
14 & โครงการพัฒนาทักษะวิชาชีพให้กับเจ้าหน้าที่ ชาวบ้าน เยาวชน พื้นที่ & 35:7 $\cdot$ 35:7 $\cdot$ 39:1 $\cdot$ 39:1 $\cdot$ 40:1 $\cdot$ 40:1\\
15 & โครงการพัฒนาต้นแบบชุดน้ำชาเซรามิกสโตนแวร์สำหรับยกระดับผลิตภัณฑ์ช & 35:8\\
16 & โครงการการถ่ายทอดองค์ความรู้กระบวนการผลิตน้ำดื่มสะอาดด้วยระบบรถน & 35:5 $\cdot$ 39:1 $\cdot$ 40:1\\
17 & โครงการพัฒนาทักษะวิชาชีพให้กับเจ้าหน้าที่ ชาวบ้าน เยาวชน พื้นที่ & 35:7 $\cdot$ 35:7 $\cdot$ 39:1 $\cdot$ 39:1 $\cdot$ 40:1 $\cdot$ 40:1\\
18 & การประยุกต์ใช้ lora mesh ชนิด multi-hops ในการส่งสัญญาณจากเซนเซอ & 35:3 $\cdot$ 39:1 $\cdot$ 40:1\\
19 & การประยุกต์ใช้เครือข่าย lora-mesh ชนิด multi-hops เพื่อการส่งสัญ & 35:3 $\cdot$ 39:1 $\cdot$ 40:1\\
20 & โครงการวิจัยความเข้มแข็งทางสังคม ชุมชนในการวางแผนการใช้น้ำศูนย์พ & 35:1\\
21 & โครงการพัฒนาทักษะอาชีพและองค์ความรู้ด้านเกษตรและเทคโนโลยีอย่างยั & 35:5 $\cdot$ 39:3 $\cdot$ 40:3\\
22 & โครงการถ่ายทอดเทคโนโลยีการเพาะขยายพันธุ์ปลาเลียหินเพื่อการอนุรัก & 35:3 $\cdot$ 39:1 $\cdot$ 40:1\\
23 & โครงส่งเสริมพัฒนาโจทย์วิจัยจากชุมชน & 35:7\\
24 & โครงการถ่ายทอดองค์ความรู้และนวัตกรรมเพื่อยกระดับมูลค่าไก่ดำ (ไก่ & 35:12 $\cdot$ 39:1 $\cdot$ 40:1\\
25 & โครงการพัฒนาระบบติดตามข้อมูลภายใต้สถานีตรวจวัดสภาพอากาศเพื่อป้อง & 35:3\\
26 & โครงการจัดทำคำขอสิ่งบ่งชี้ทางภูมิศาสตร์ไทย (GI) สินค้ากาแฟเลอตอ & 38:1 $\cdot$ 39:1 $\cdot$ 39:1 $\cdot$ 40:1 $\cdot$ 40:1\\
\midrule
\bottomrule
\end{longtable}
}


\clearpage
% ================= ขอบคุณ =================
\section*{กิตติกรรมประกาศ}
\addcontentsline{toc}{section}{กิตติกรรมประกาศ}

การดำเนินงานของกลุ่มแผนงานใต้ร่มพระบารมี มหาวิทยาลัยเทคโนโลยีราชมงคลล้านนา ปีงบประมาณ 2569
สำเร็จลุล่วงมาถึงรอบ 9 เดือนได้ด้วยความร่วมมือจากหน่วยงานและบุคคลต่อไปนี้

\begin{center}\small
\setlength{\tabcolsep}{8pt}\renewcommand{\arraystretch}{1.3}
\begin{tabular}{p{4.7cm}p{4.7cm}p{4.7cm}}
\toprule
\rowcolor{rpfnavy!12}
\textbf{หน่วยงานในพื้นที่} & \textbf{หน่วยงานภายในมหาวิทยาลัย} & \textbf{วิทยาเขตและเครือข่าย}\\
\midrule
มูลนิธิโครงการหลวง ;
ศูนย์พัฒนาโครงการหลวงและสถานีเกษตรหลวงในพื้นที่ ;
ศูนย์พัฒนาพันธุ์พืชจักรพันธ์เพ็ญศิริ ;
ศูนย์ศึกษาการพัฒนาห้วยฮ่องไคร้อันเนื่องมาจากพระราชดำริ
&
กลุ่มแผนงานใต้ร่มพระบารมี ;
สถาบันถ่ายทอดเทคโนโลยีสู่ชุมชน ;
สถาบันวิจัยเทคโนโลยีเกษตร ;
คณะวิศวกรรมศาสตร์ ;
คณะศิลปกรรมและสถาปัตยกรรมศาสตร์ ;
คณะบริหารธุรกิจและศิลปศาสตร์ ;
วิทยาลัยเทคโนโลยีและสหวิทยาการ
&
มทร.ล้านนา เชียงราย น่าน ตาก และพิษณุโลก ;
โรงเรียนพระปริยัติธรรมและสถานศึกษาในพื้นที่ ;
วัดและวิสาหกิจชุมชนร่วมโครงการ ;
สำนักงานเขตพื้นที่การศึกษาประถมศึกษาแม่ฮ่องสอน เขต 2
\\
\bottomrule
\end{tabular}
\end{center}

ขอขอบคุณหัวหน้าโครงการทั้ง 42 ท่าน และทีมงานทุกท่านที่ร่วมขับเคลื่อนงานในพื้นที่
รวมถึงเจ้าหน้าที่ผู้ดูแลข้อมูลงบประมาณและระบบติดตามโครงการที่ทำให้รายงานฉบับนี้ตั้งอยู่บนข้อมูลจริง

\vspace{14pt}
{\footnotesize\color{rpfink}\textbf{จัดทำโดย} กลุ่มแผนงานใต้ร่มพระบารมี มหาวิทยาลัยเทคโนโลยีราชมงคลล้านนา\par
\textbf{ข้อมูล ณ} 22 กรกฎาคม 2569 $\cdot$ รายงานฉบับสมบูรณ์จะจัดทำเมื่อสิ้นปีงบประมาณ 30 กันยายน 2569}

\end{document}
